\documentclass{amsart}
\usepackage{amsmath,amssymb,amsthm,hyperref}

\newtheorem{theorem}{Theorem}
\newtheorem{lemma}{Lemma}
\newtheorem{proposition}{Proposition}
\newtheorem{corollary}{Corollary}
\theoremstyle{definition}
\newtheorem{definition}{Definition}
\theoremstyle{remark}
\newtheorem{remark}{Remark}

\newcommand{\R}{\mathbb{R}}
\newcommand{\E}{\mathbb{E}}
\newcommand{\norm}[1]{\left\lVert#1\right\rVert}
\DeclareMathOperator{\diag}{diag}

\begin{document}

\title{Existence and pathwise uniqueness for the queueing-network reflected Brownian motion}
\maketitle

\section{Formalisation and statement of the result}

Throughout, vectors are columns in $\R^d$, inequalities between vectors or
matrices are understood entrywise, $I$ is the $d\times d$ identity, and for a
square matrix $A$ we write $\rho(A)$ for its spectral radius. For a function
$f:[0,\infty)\to\R^d$ and $T\ge 0$ we write
\[
  \norm{f}_T:=\sup_{0\le t\le T}\max_{1\le j\le d}|f_j(t)|.
\]
We denote by $C_0=C_0([0,\infty);\R^d)$ the set of continuous paths
$x:[0,\infty)\to\R^d$ with $x(0)=0$.

We are given $G\in\R^{d\times d}$, $\theta\in\R^d$, a symmetric positive
semidefinite $\Gamma\in\R^{d\times d}$, and a $d$-dimensional Brownian motion
$\xi^\ast$ with drift $-\theta$ and covariance $\Gamma$, $\xi^\ast(0)=0$, on a
filtered probability space $(\Omega,\mathcal F,(\mathcal F_t),\mathbb P)$
satisfying the usual conditions. We seek continuous adapted
$W^\ast,Y^\ast:[0,\infty)\times\Omega\to\R^d$ such that, almost surely, for all
$t\ge 0$ and all $j$:
\begin{align}
  W^\ast(t)&=\xi^\ast(t)-G\,W^\ast(t)+Y^\ast(t),\label{eq:dyn}\\
  W^\ast_j(t)&\ge 0,\label{eq:pos}\\
  Y^\ast_j(0)&=0,\quad Y^\ast_j\text{ continuous and nondecreasing},\label{eq:reg}\\
  \int_0^\infty &W^\ast_j(t)\,dY^\ast_j(t)=0.\label{eq:comp}
\end{align}

\paragraph{Reduction to a Skorokhod problem.}
Set
\begin{equation}\label{eq:RQ}
  R:=I+G .
\end{equation}
Equation \eqref{eq:dyn} is equivalent to $R\,W^\ast(t)=\xi^\ast(t)+Y^\ast(t)$.
Assume for the moment that $R$ is invertible (this is part of the hypothesis of
Theorem~\ref{thm:main}). Then \eqref{eq:dyn} is equivalent to
\begin{equation}\label{eq:SPform}
  W^\ast(t)=X^\ast(t)+R^{-1}Y^\ast(t),\qquad
  X^\ast(t):=R^{-1}\xi^\ast(t).
\end{equation}
Writing
\begin{equation}\label{eq:Qdef}
  Q:=I-R^{-1}=G(I+G)^{-1}=R^{-1}G,
\end{equation}
(the last two equalities hold because $R^{-1}=(I+G)^{-1}$ commutes with $G$ and
$I-(I+G)^{-1}=(I+G-I)(I+G)^{-1}=G(I+G)^{-1}$), we may rewrite
\eqref{eq:SPform} as
\begin{equation}\label{eq:SPstd}
  W^\ast(t)=X^\ast(t)+Y^\ast(t)-Q\,Y^\ast(t).
\end{equation}
The problem \eqref{eq:dyn}--\eqref{eq:comp} is thus, pathwise, an instance of
the \emph{Skorokhod problem} on the nonnegative orthant $\R_+^d$ with input
$X^\ast$ and reflection matrix $R^{-1}=I-Q$: given $x\in C_0$, find continuous
$z,y:[0,\infty)\to\R^d$ with
\begin{align}
  z(t)&=x(t)+y(t)-Q\,y(t),\label{eq:sp1}\\
  z_j(t)&\ge 0,\qquad j=1,\dots,d,\label{eq:sp2}\\
  y_j(0)&=0,\quad y_j\text{ continuous nondecreasing},\label{eq:sp3}\\
  \int_0^\infty z_j&(t)\,dy_j(t)=0,\qquad j=1,\dots,d.\label{eq:sp4}
\end{align}
A pair $(z,y)$ satisfying \eqref{eq:sp1}--\eqref{eq:sp4} is called a
\emph{solution of the Skorokhod problem $\mathrm{SP}(x;Q)$}.

Our main result is the following. It identifies an explicit, verifiable class of
triples for which existence and pathwise uniqueness hold, shows that this class
contains every triple arising from a multiclass queueing network, and explains
precisely why no such positive statement can hold for an arbitrary triple
$(\theta,\Gamma,G)$.

\begin{theorem}\label{thm:main}
Let $G\in\R^{d\times d}$ and put $R=I+G$, $Q=G(I+G)^{-1}$ as in \eqref{eq:RQ},
\eqref{eq:Qdef}. Suppose
\begin{equation}\label{eq:HR}
  R=I+G\ \text{is invertible},\qquad Q\ge 0\ \text{(entrywise)},\qquad \rho(Q)<1 .
  \tag{HR}
\end{equation}
Then for every $\theta\in\R^d$ and every symmetric positive semidefinite
$\Gamma$ there exist continuous adapted processes $(W^\ast,Y^\ast)$ satisfying
\eqref{eq:dyn}--\eqref{eq:comp} almost surely, and this pair is pathwise unique:
if $(W^\ast,Y^\ast)$ and $(\widetilde W^\ast,\widetilde Y^\ast)$ are two such
pairs then $\mathbb P\big(W^\ast(t)=\widetilde W^\ast(t),\,
Y^\ast(t)=\widetilde Y^\ast(t)\ \forall t\ge 0\big)=1$.
Moreover $W^\ast$ is a $(\mathcal F_t)$-semimartingale reflecting Brownian
motion, and $W^\ast,Y^\ast$ are obtained from the sample paths of $\xi^\ast$ by
a fixed deterministic Lipschitz map.
\end{theorem}

\section{The deterministic Skorokhod map}

The heart of the matter is the following purely analytic statement
(Harrison--Reiman). Condition \eqref{eq:HR} is imposed throughout this section.

\begin{theorem}\label{thm:HR}
Assume \eqref{eq:HR}. For every $x\in C_0$ there exists a unique pair
$(z,y)\in C_0\times C_0$ solving $\mathrm{SP}(x;Q)$. Writing
$z=\Phi(x)$, $y=\Psi(x)$, the maps $\Phi,\Psi:C_0\to C_0$ are Lipschitz on
every interval: there is a constant $L=L(Q)<\infty$ such that for all
$x,\tilde x\in C_0$ and all $T\ge 0$,
\begin{equation}\label{eq:lip}
  \norm{\Phi(x)-\Phi(\tilde x)}_T+\norm{\Psi(x)-\Psi(\tilde x)}_T
  \le L\,\norm{x-\tilde x}_T .
\end{equation}
\end{theorem}

The proof uses the elementary one--dimensional Skorokhod map and a fixed-point
argument; the spectral assumption $\rho(Q)<1$ together with $Q\ge0$ enters only
through the Perron--Frobenius lemma below.

\subsection{The one-dimensional Skorokhod map}

\begin{lemma}\label{lem:1d}
Let $u\in C([0,\infty);\R)$ with $u(0)\ge 0$. Define
\begin{equation}\label{eq:1dmap}
  \ell(t):=\max_{0\le s\le t}\big(-u(s)\big)^+=\Big(\,{-}\!\min_{0\le s\le t}u(s)\Big)^+,
  \qquad p(t):=u(t)+\ell(t).
\end{equation}
Then $(p,\ell)$ is the unique pair of continuous functions such that
$p=u+\ell$, $p\ge0$ on $[0,\infty)$, $\ell(0)=0$, $\ell$ is nondecreasing, and
$\int_0^\infty p\,d\ell=0$. Moreover the maps $u\mapsto\ell$ and $u\mapsto p$
are Lipschitz in the sup norm: for $u,\tilde u$ and every $T$,
\begin{equation}\label{eq:1dlip}
  \sup_{0\le t\le T}|\ell(t)-\tilde\ell(t)|\le \sup_{0\le t\le T}|u(t)-\tilde u(t)|,
  \qquad
  \sup_{0\le t\le T}|p(t)-\tilde p(t)|\le 2\sup_{0\le t\le T}|u(t)-\tilde u(t)|.
\end{equation}
Finally, $\ell$ is monotone in $u$: if $u\le\tilde u$ on $[0,T]$ then
$\ell\ge\tilde\ell$ on $[0,T]$.
\end{lemma}

\begin{proof}
\emph{Existence and the stated properties.} Since $u(0)\ge0$ we have
$\ell(0)=(\,{-}u(0))^+=0$. The function $t\mapsto\min_{0\le s\le t}u(s)$ is
continuous and nonincreasing, hence $\ell$ is continuous and nondecreasing, and
$p=u+\ell$ is continuous. For each $t$, $\ell(t)\ge -u(t)$, so
$p(t)=u(t)+\ell(t)\ge0$, giving \eqref{eq:sp2}. For the complementarity
condition, $\ell$ increases only at times $t$ where
$\min_{0\le s\le t}u(s)=u(t)<0$ (a point of new running minimum with negative
value); at such $t$, $\ell(t)=-u(t)$, so $p(t)=0$. Hence $\{t:p(t)>0\}$ is
contained in the set on which $\ell$ is locally constant, and the Stieltjes
integral $\int_0^\infty p\,d\ell=0$.

\emph{Uniqueness.} Suppose $(p,\ell)$ and $(p',\ell')$ both satisfy the four
conditions. Put $h:=\ell-\ell'=p-p'$ (the two expressions are equal since
$p-p'=(u+\ell)-(u+\ell')$). Fix $t$ and consider
$\tfrac12\frac{d}{ds}h(s)^2=h(s)\,dh(s)=(\ell-\ell')(s)\,(d\ell-d\ell')(s)$ in
the Stieltjes sense. Now
$\int_0^t(\ell-\ell')\,d\ell=\int_0^t(p-p')\,d\ell
=\int_0^t p\,d\ell-\int_0^t p'\,d\ell=-\int_0^t p'\,d\ell\le 0$,
because $\int p\,d\ell=0$ and $p'\ge0$, $d\ell\ge0$. Symmetrically
$\int_0^t(\ell-\ell')\,(-d\ell')=\int_0^t(p'-p)\,d\ell'
=-\int_0^t p\,d\ell'\le0$. Hence
$\tfrac12 h(t)^2=\int_0^t h\,dh
=\int_0^t(\ell-\ell')d\ell-\int_0^t(\ell-\ell')d\ell'\le 0$, so $h\equiv0$,
i.e.\ $\ell=\ell'$ and $p=p'$. (Here we used $h(0)=0$ and that $h$ is of
bounded variation and continuous, so the chain rule
$\tfrac12 h(t)^2=\int_0^t h\,dh$ is valid.)

\emph{Lipschitz and monotonicity bounds.} For the first inequality in
\eqref{eq:1dlip}, fix $t\le T$ and let $s_0\le t$ attain
$\min_{0\le s\le t}u(s)$. Then
$\ell(t)=(-u(s_0))^+$ and
$\tilde\ell(t)=\max_{s\le t}(-\tilde u(s))^+\ge(-\tilde u(s_0))^+
\ge -\tilde u(s_0)\ge -u(s_0)-\norm{u-\tilde u}_T$, while if $\ell(t)=0$ the
inequality $\ell(t)-\tilde\ell(t)\le \norm{u-\tilde u}_T$ is trivial; in either
case $\ell(t)-\tilde\ell(t)\le \norm{u-\tilde u}_T$. By symmetry
$|\ell(t)-\tilde\ell(t)|\le\norm{u-\tilde u}_T$. The second inequality follows
from $p-\tilde p=(u-\tilde u)+(\ell-\tilde\ell)$ and the triangle inequality.
Monotonicity is immediate from \eqref{eq:1dmap}: if $u\le\tilde u$ then
$-u\ge-\tilde u$ pointwise, so $\ell\ge\tilde\ell$.
\end{proof}

\subsection{A weighted norm from Perron--Frobenius}

\begin{lemma}\label{lem:weight}
Assume $Q\ge0$ and $\rho:=\rho(Q)<1$. Then there exist a strictly positive
vector $w\in\R^d$, $w>0$, and a constant $\gamma\in[\rho,1)$ such that
\begin{equation}\label{eq:Qw}
  Q\,w\le \gamma\,w \quad\text{(entrywise).}
\end{equation}
Consequently, defining the weighted sup norm
$\norm{v}_w:=\max_{1\le j\le d}|v_j|/w_j$ on $\R^d$, the linear map
$v\mapsto Qv$ satisfies, for every $v\ge0$,
$\norm{Qv}_w\le\gamma\,\norm{v}_w$ (only this nonnegative case is used
below).
\end{lemma}

\begin{proof}
Replace $Q$ by $Q_\varepsilon:=Q+\varepsilon\mathbf 1\mathbf 1^\top$ with
$\varepsilon>0$ small, where $\mathbf 1$ is the all-ones vector; $Q_\varepsilon$
is strictly positive (all entries $>0$) and $\rho(Q_\varepsilon)\to\rho(Q)<1$ as
$\varepsilon\downarrow0$ (the spectral radius is continuous in the matrix
entries), so we may fix $\varepsilon$ with $\gamma:=\rho(Q_\varepsilon)<1$. By
the Perron--Frobenius theorem applied to the strictly positive matrix
$Q_\varepsilon$, there is a strictly positive eigenvector $w>0$ with
$Q_\varepsilon w=\gamma w$. Since $Q\le Q_\varepsilon$ entrywise and $w>0$, we
get $Qw\le Q_\varepsilon w=\gamma w$, which is \eqref{eq:Qw} with this $w$ and
$\gamma\in(\rho,1)$ (and $\gamma\ge\rho$ because $\rho$ is monotone under
entrywise domination of nonnegative matrices).

For $v\ge0$, $0\le Qv\le Q\big(\max_k(v_k/w_k)\,w\big)
=\norm{v}_w\,Qw\le \gamma\norm{v}_w\,w$, whence
$(Qv)_j/w_j\le\gamma\norm{v}_w$ for every $j$, i.e.\
$\norm{Qv}_w\le\gamma\norm{v}_w$. (We only use \eqref{eq:Qw} for $v\ge0$ in the
sequel.)
\end{proof}

We also fix the equivalence constants of the weighted norm: with
$\underline w:=\min_j w_j>0$ and $\overline w:=\max_j w_j$,
\begin{equation}\label{eq:equiv}
  \underline w\,\norm{v}_w\le \max_j|v_j|\le \overline w\,\norm{v}_w
  \qquad(v\in\R^d).
\end{equation}
For path-valued functions we use the corresponding norm
$\norm{f}_{w,T}:=\sup_{0\le t\le T}\norm{f(t)}_w$.

\subsection{Proof of Theorem~\ref{thm:HR}}

Fix $x\in C_0$ and $T>0$. We seek $y\in C([0,T];\R^d)$, $y$ nondecreasing
componentwise with $y(0)=0$, such that with
\begin{equation}\label{eq:fixedpt}
  z=x+y-Qy,
\end{equation}
the pair $(z,y)$ solves $\mathrm{SP}(x;Q)$ on $[0,T]$.

\paragraph{Coordinatewise decoupling.}
For each coordinate $j$, conditions \eqref{eq:sp2}--\eqref{eq:sp4} for that
coordinate say exactly that $(z_j,y_j)$ is the one-dimensional Skorokhod
reflection of the input
\begin{equation}\label{eq:input}
  u_j:=x_j-(Qy)_j,
\end{equation}
because from \eqref{eq:fixedpt}, $z_j=u_j+y_j$, and \eqref{eq:input} has
$u_j(0)=x_j(0)-(Qy(0))_j=0\ge0$. By Lemma~\ref{lem:1d}, given $y$ the unique
continuous $z_j\ge0$, nondecreasing $y_j$ with $y_j(0)=0$ and
$\int z_j\,dy_j=0$ satisfying $z_j=u_j+y_j$ are
\begin{equation}\label{eq:Tdef}
  y_j=\Lambda_j(y):=\Big({-}\!\min_{0\le s\le\cdot}u_j(s)\Big)^+
  =\Big({-}\!\min_{0\le s\le\cdot}\big(x_j-(Qy)_j\big)(s)\Big)^+ .
\end{equation}
Thus $(z,y)$ solves $\mathrm{SP}(x;Q)$ on $[0,T]$ if and only if $y$ is a fixed
point of the map $\Lambda=(\Lambda_1,\dots,\Lambda_d)$ acting on
\[
  \mathcal K_T:=\{y\in C([0,T];\R^d):y(0)=0,\ y_j\text{ nondecreasing}\}.
\]

\paragraph{$\Lambda$ maps $\mathcal K_T$ into itself.}
For $y\in\mathcal K_T$ each $u_j=x_j-(Qy)_j$ is continuous with
$u_j(0)=0\ge0$, so by Lemma~\ref{lem:1d} each $\Lambda_j(y)$ is continuous,
nondecreasing, with $\Lambda_j(y)(0)=0$. Hence $\Lambda:\mathcal K_T\to\mathcal
K_T$.

\paragraph{$\Lambda$ is a contraction in the weighted norm.}
Let $y,\tilde y\in\mathcal K_T$ and put $u_j=x_j-(Qy)_j$,
$\tilde u_j=x_j-(Q\tilde y)_j$, so that
$u_j-\tilde u_j=-(Q(y-\tilde y))_j$. By the first inequality in
\eqref{eq:1dlip} of Lemma~\ref{lem:1d}, applied on $[0,t]$ for $t\le T$ and
coordinate $j$,
\begin{equation}\label{eq:perj}
  \sup_{0\le s\le t}\big|\Lambda_j(y)(s)-\Lambda_j(\tilde y)(s)\big|
  \le \sup_{0\le s\le t}\big|u_j(s)-\tilde u_j(s)\big|
  = \sup_{0\le s\le t}\big|\big(Q(y-\tilde y)\big)_j(s)\big|.
\end{equation}
Fix $t\le T$ and write $\delta:=y-\tilde y$. For each fixed $s\le t$ we estimate
$|(Q\delta)_j(s)|$. Since $Q\ge0$,
\begin{equation}\label{eq:absbound}
  \big|(Q\delta(s))_j\big|=\Big|\sum_k Q_{jk}\delta_k(s)\Big|
  \le \sum_k Q_{jk}\,|\delta_k(s)|=\big(Q\,|\delta(s)|\big)_j ,
\end{equation}
where $|\delta(s)|:=(|\delta_1(s)|,\dots,|\delta_d(s)|)^\top\ge0$. By
Lemma~\ref{lem:weight} (the nonnegative case), for the vector
$a(s):=|\delta(s)|\ge0$ we have
$\norm{Q a(s)}_w\le\gamma\norm{a(s)}_w$, i.e.\
$(Q|\delta(s)|)_j/w_j\le\gamma\,\norm{\delta(s)}_w$ for every $j$. Combining
with \eqref{eq:absbound},
\begin{equation}\label{eq:weightj}
  \frac{|(Q\delta(s))_j|}{w_j}\le\gamma\,\norm{\delta(s)}_w
  \le\gamma\,\norm{\delta}_{w,t},\qquad s\le t .
\end{equation}
Dividing \eqref{eq:perj} by $w_j$ and using \eqref{eq:weightj},
\[
  \frac1{w_j}\sup_{0\le s\le t}\big|\Lambda_j(y)(s)-\Lambda_j(\tilde y)(s)\big|
  \le \frac1{w_j}\sup_{0\le s\le t}\big|(Q\delta(s))_j\big|
  \le \gamma\,\norm{\delta}_{w,t}.
\]
Taking the maximum over $j$ and the supremum over $s\le t$ gives
\begin{equation}\label{eq:contraction}
  \norm{\Lambda(y)-\Lambda(\tilde y)}_{w,t}\le\gamma\,\norm{y-\tilde y}_{w,t},
  \qquad 0\le t\le T,
\end{equation}
with $\gamma\in[\rho(Q),1)$ independent of $t,T,x$.

\paragraph{Existence and uniqueness of the fixed point.}
$\big(\mathcal K_T,\norm{\cdot}_{w,T}\big)$ is a closed subset of the Banach
space $\big(C([0,T];\R^d),\norm{\cdot}_{w,T}\big)$ (uniform limits of
nondecreasing functions with value $0$ at $0$ are again such), hence a complete
metric space. By \eqref{eq:contraction} with $t=T$, $\Lambda$ is a
$\gamma$-contraction on it, $\gamma<1$. The Banach fixed-point theorem yields a
unique $y\in\mathcal K_T$ with $\Lambda(y)=y$; setting $z:=x+y-Qy$ produces the
(unique, by the equivalence above) solution $(z,y)$ of $\mathrm{SP}(x;Q)$ on
$[0,T]$.

\paragraph{Consistency in $T$ and global solution.}
If $T'<T$, the restriction to $[0,T']$ of the solution on $[0,T]$ solves
$\mathrm{SP}(x;Q)$ on $[0,T']$ (all four conditions are local in $t$), so by
uniqueness it coincides with the solution constructed on $[0,T']$. Hence the
solutions are consistent, and define a unique global solution
$(z,y)=(\Phi(x),\Psi(x))$ on $[0,\infty)$.

\paragraph{Lipschitz dependence \eqref{eq:lip}.}
Let $x,\tilde x\in C_0$ with fixed points $y=\Psi(x)$,
$\tilde y=\Psi(\tilde x)$ on $[0,T]$. Repeating \eqref{eq:perj} with the two
\emph{different} inputs $u_j=x_j-(Qy)_j$ and $\tilde u_j=\tilde
x_j-(Q\tilde y)_j$, now $u_j-\tilde u_j=(x_j-\tilde x_j)-(Q(y-\tilde y))_j$, and
Lemma~\ref{lem:1d} gives, for $t\le T$ and every $j$,
\[
  \sup_{s\le t}|y_j(s)-\tilde y_j(s)|
  \le \sup_{s\le t}|x_j(s)-\tilde x_j(s)|+\sup_{s\le t}|(Q(y-\tilde y))_j(s)|.
\]
Dividing by $w_j$, using \eqref{eq:weightj} for the second term and
$\sup_{s\le t}|x_j(s)-\tilde x_j(s)|/w_j\le (1/\underline w)\norm{x-\tilde x}_t$
(from $w_j\ge\underline w$), and maximising over $j$,
\[
  \norm{y-\tilde y}_{w,t}\le \frac1{\underline w}\,\norm{x-\tilde x}_{t}
  +\gamma\,\norm{y-\tilde y}_{w,t}.
\]
Since $\gamma<1$, $\norm{y-\tilde y}_{w,T}\le
\frac{1}{\underline w(1-\gamma)}\,\norm{x-\tilde x}_{T}$, and by
\eqref{eq:equiv},
\begin{equation}\label{eq:psilip}
  \norm{\Psi(x)-\Psi(\tilde x)}_T\le \overline w\,\norm{y-\tilde y}_{w,T}
  \le \frac{\overline w}{\underline w(1-\gamma)}\,\norm{x-\tilde x}_T=:L_1\norm{x-\tilde x}_T .
\end{equation}
For $\Phi$, $z-\tilde z=(x-\tilde x)+(I-Q)(y-\tilde y)$, so with
$\norm{(I-Q)}_\infty:=\max_j\sum_k|(I-Q)_{jk}|$ (the matrix $\ell^\infty$
operator norm),
\begin{equation}\label{eq:philip}
  \norm{\Phi(x)-\Phi(\tilde x)}_T\le\norm{x-\tilde x}_T
  +\norm{I-Q}_\infty\,\norm{\Psi(x)-\Psi(\tilde x)}_T
  \le\big(1+\norm{I-Q}_\infty L_1\big)\norm{x-\tilde x}_T .
\end{equation}
Adding \eqref{eq:psilip} and \eqref{eq:philip} gives \eqref{eq:lip} with
$L:=1+L_1+\norm{I-Q}_\infty L_1$. This proves Theorem~\ref{thm:HR}.
\qed

\section{Proof of Theorem~\ref{thm:main}}

\paragraph{Construction.}
Define $X^\ast(t):=R^{-1}\xi^\ast(t)=(I+G)^{-1}\xi^\ast(t)$. Since $\xi^\ast$
has continuous sample paths with $\xi^\ast(0)=0$, almost every path of
$X^\ast$ lies in $C_0$. Set
\begin{equation}\label{eq:WYdef}
  W^\ast:=\Phi(X^\ast),\qquad Y^\ast:=\Psi(X^\ast),
\end{equation}
the deterministic Skorokhod maps of Theorem~\ref{thm:HR} applied pathwise.

\paragraph{$(W^\ast,Y^\ast)$ satisfies \eqref{eq:dyn}--\eqref{eq:comp}.}
By Theorem~\ref{thm:HR}, almost surely $(W^\ast,Y^\ast)=(z,y)$ solves
$\mathrm{SP}(X^\ast;Q)$, i.e.\ \eqref{eq:sp1}--\eqref{eq:sp4} hold with
$x=X^\ast$. Conditions \eqref{eq:sp2}, \eqref{eq:sp3}, \eqref{eq:sp4} are
exactly \eqref{eq:pos}, \eqref{eq:reg}, \eqref{eq:comp}. For \eqref{eq:dyn}, by
\eqref{eq:sp1} and \eqref{eq:Qdef},
\[
  W^\ast=X^\ast+(I-Q)Y^\ast=R^{-1}\xi^\ast+R^{-1}Y^\ast,
\]
hence $R\,W^\ast=\xi^\ast+Y^\ast$, that is
$(I+G)W^\ast=\xi^\ast+Y^\ast$, equivalently
$W^\ast=\xi^\ast-G\,W^\ast+Y^\ast$, which is \eqref{eq:dyn}.

\paragraph{Continuity, adaptedness, measurability.}
$W^\ast,Y^\ast$ are continuous because $\Phi(x),\Psi(x)\in C_0$ for every
$x\in C_0$. For adaptedness, note from \eqref{eq:Tdef} and \eqref{eq:fixedpt}
that the values $W^\ast(t),Y^\ast(t)$ depend only on $\{X^\ast(s):s\le t\}$,
hence only on $\{\xi^\ast(s):s\le t\}$. Concretely, the Picard iterates
$y^{(0)}\equiv0$, $y^{(n+1)}=\Lambda(y^{(n)})$ converge uniformly on $[0,t]$ to
$Y^\ast$ by \eqref{eq:contraction}; each $y^{(n)}(t)$ is, by \eqref{eq:Tdef}, a
continuous functional of $\{X^\ast(s):s\le t\}$ and hence
$\mathcal F_t$-measurable (the running minimum and the maps in
Lemma~\ref{lem:1d} are continuous in the sup norm, by \eqref{eq:1dlip}). The
uniform limit $Y^\ast(t)$ is therefore $\mathcal F_t$-measurable, and so is
$W^\ast(t)=X^\ast(t)+(I-Q)Y^\ast(t)$. Thus $(W^\ast,Y^\ast)$ is adapted.

\paragraph{Pathwise uniqueness.}
Suppose $(W^\ast,Y^\ast)$ and $(\widetilde W^\ast,\widetilde Y^\ast)$ both
satisfy \eqref{eq:dyn}--\eqref{eq:comp} a.s. On the full-probability event where
both hold and $\xi^\ast$ is continuous, each pair solves
$\mathrm{SP}(X^\ast;Q)$ pathwise (by the equivalence \eqref{eq:SPstd}). By the
uniqueness in Theorem~\ref{thm:HR}, $W^\ast=\widetilde W^\ast$ and
$Y^\ast=\widetilde Y^\ast$ as functions on $[0,\infty)$ on that event, which has
probability one.

\paragraph{$W^\ast$ is a semimartingale reflecting Brownian motion.}
$\xi^\ast(t)=B(t)-\theta t$ with $B$ a continuous $(\mathcal F_t)$-martingale
(Brownian motion with covariance $\Gamma$), so $X^\ast=R^{-1}\xi^\ast$ is a
continuous semimartingale with martingale part $R^{-1}B$ and finite-variation
part $-R^{-1}\theta\,t$. By \eqref{eq:WYdef}, \eqref{eq:sp1},
\[
  W^\ast=R^{-1}\xi^\ast+(I-Q)Y^\ast
       = R^{-1}B(t)-R^{-1}\theta\,t+R^{-1}Y^\ast,
\]
since $(I-Q)=R^{-1}$. The right-hand side is a continuous adapted process equal
to a continuous local martingale ($R^{-1}B$) plus a continuous adapted
finite-variation process ($-R^{-1}\theta\,t+R^{-1}Y^\ast$, the latter because
each $Y^\ast_j$ is nondecreasing). Hence $W^\ast$ is a continuous
$(\mathcal F_t)$-semimartingale taking values in $\R_+^d$, with constraining
term $R^{-1}Y^\ast=(I-Q)Y^\ast$, drift $-R^{-1}\theta=-(I-Q)\theta$ and
covariance $R^{-1}\Gamma R^{-\top}=(I-Q)\Gamma(I-Q)^\top$; this is precisely a
semimartingale reflecting Brownian motion in $\R_+^d$ with reflection matrix
$R^{-1}=I-Q$. This completes the proof of Theorem~\ref{thm:main}.
\qed

\section{Multiclass queueing networks, and the general triple}

\subsection{Every queueing-network triple satisfies \eqref{eq:HR}}

In the heavy-traffic (diffusion) approximation of an open multiclass queueing
network, the reflection matrix of the limiting workload/queue-length process is
\begin{equation}\label{eq:queueR}
  R^{-1}=I-P^\top,
\end{equation}
where $P=(P_{kj})$ is the (substochastic) routing matrix of the network:
$P_{kj}\ge0$ is the probability that a job leaving class $k$ goes next to class
$j$, $\sum_j P_{kj}\le1$, and the network is open, meaning every job eventually
leaves the system; equivalently $P$ is transient, $\rho(P)<1$, and
$\sum_{n\ge0}P^n=(I-P)^{-1}$ is finite. Comparing \eqref{eq:queueR} with
\eqref{eq:Qdef} identifies
\[
  Q=P^\top, \qquad R=I+G=(I-P^\top)^{-1},\qquad
  G=(I-P^\top)^{-1}-I=P^\top(I-P^\top)^{-1}\ge0 .
\]
Then \eqref{eq:HR} holds automatically: $R=(I-P^\top)^{-1}$ is invertible;
$Q=P^\top\ge0$ entrywise; and $\rho(Q)=\rho(P^\top)=\rho(P)<1$ by openness of
the network. (Here $G$ is the matrix whose $(j,k)$ entry equals the expected
number of future visits to class $j$ by a job currently at class $k$, the
``routing-information'' interpretation.) Therefore
Theorem~\ref{thm:main} applies, and \emph{existence and pathwise uniqueness of
$(W^\ast,Y^\ast)$ hold for every triple $(\theta,\Gamma,G)$ arising from an open
multiclass queueing network}, with no restriction on $\theta$ or on the
positive semidefinite $\Gamma$.

\subsection{An arbitrary triple: failure of \eqref{eq:HR} and what survives}

For a completely arbitrary $G$ the answer is negative: existence (let alone
uniqueness) of a pathwise solution can fail. Two phenomena occur.

\emph{(a) $R=I+G$ singular.} If $-1$ is an eigenvalue of $G$, then $R=I+G$ is
singular and the linear relation $(I+G)W^\ast=\xi^\ast+Y^\ast$ generically has
no solution: choosing $\theta=0$ and $\Gamma$ nondegenerate, $\xi^\ast(t)$ a.s.\
leaves the column space of $R$ shifted by $-\mathrm{range}\,Y^\ast$ at small
times, so \eqref{eq:dyn} cannot hold. Hence invertibility of $I+G$ is necessary
for \eqref{eq:dyn} to be solvable for all directions of the noise.

\emph{(b) $R$ invertible but $Q=G(I+G)^{-1}\not\ge0$ or $\rho(Q)\ge1$.} The
existence theory for the Skorokhod problem on $\R_+^d$ for \emph{all} continuous
inputs requires the reflection matrix $R^{-1}=I-Q$ to be \emph{completely-}$\mathcal S$,
i.e.\ every principal submatrix of $R^{-1}$ is an $\mathcal S$-matrix (Bernard--El
Kharroubi; Taylor--Williams): a $k\times k$ matrix $M$ is an $\mathcal S$-matrix
if there exists $\lambda\ge0$ with $M\lambda>0$ entrywise. When $R^{-1}=I-Q$ is
not completely-$\mathcal S$ there exist continuous inputs $x$ (and, for
nondegenerate $\Gamma$, the Brownian path $X^\ast$ with positive probability)
for which no solution $(z,y)$ of $\mathrm{SP}(x;Q)$ exists; consequently no
$(W^\ast,Y^\ast)$ satisfying \eqref{eq:dyn}--\eqref{eq:comp} exists.
Condition \eqref{eq:HR} ($Q\ge0$, $\rho(Q)<1$) is a clean sufficient condition
that implies $R^{-1}=I-Q$ is a nonsingular $M$-matrix --- indeed
$(I-Q)^{-1}=\sum_{n\ge0}Q^n\ge0$ converges since $\rho(Q)<1$ --- and every
nonsingular $M$-matrix is completely-$\mathcal S$ (take $\lambda=(I-Q)^{-1}\mathbf
1>0$, and the same for each principal submatrix, which is again a nonsingular
$M$-matrix). Thus \eqref{eq:HR} guarantees not merely existence but the strong,
pathwise, Lipschitz solution of Theorem~\ref{thm:main}.

\emph{(c) Existence without pathwise uniqueness.} When $R^{-1}$ is
completely-$\mathcal S$ but \eqref{eq:HR} fails (so the deterministic Skorokhod
map need not be single-valued/Lipschitz), one can still construct a
\emph{semimartingale reflecting Brownian motion} $W^\ast$ as a weak solution
that is \emph{unique in law}, by the submartingale-problem method of
Reiman--Williams and Taylor--Williams, provided $R^{-1}$ is completely-$\mathcal S$;
but the pair $(W^\ast,Y^\ast)$ need not be pathwise unique and need not be a
measurable functional of $\xi^\ast$'s path. Pathwise uniqueness in the sense of
Theorem~\ref{thm:main} is exactly what \eqref{eq:HR} delivers.

\subsection{Summary answer}

\begin{itemize}
\item There exists a pathwise solution $(W^\ast,Y^\ast)$ of
\eqref{eq:dyn}--\eqref{eq:comp}, for every drift $\theta$ and every PSD
$\Gamma$, and it is pathwise unique and a Lipschitz functional of the path of
$\xi^\ast$, \textbf{whenever} $I+G$ is invertible and $Q=G(I+G)^{-1}$ is
nonnegative with $\rho(Q)<1$ (the Harrison--Reiman condition \eqref{eq:HR}).
\item This class of triples \textbf{contains every triple arising from an open
multiclass queueing network}, because there $G=P^\top(I-P^\top)^{-1}$ with $P$
the (transient, substochastic) routing matrix, so $Q=P^\top\ge0$ and
$\rho(Q)=\rho(P)<1$. Hence the answer to the final question is: \emph{yes},
existence and pathwise uniqueness hold for every queueing-network triple.
\item For a \textbf{general} triple the positive statement fails: invertibility
of $I+G$ is necessary for \eqref{eq:dyn} to be solvable in all noise directions,
and existence of solutions for all continuous (in particular all Brownian)
inputs requires $I-Q$ to be completely-$\mathcal S$. Beyond \eqref{eq:HR} one
generally obtains at best weak existence and uniqueness in law (under
completely-$\mathcal S$), not pathwise uniqueness.
\end{itemize}

\end{document}
